% ---
% id: EX-CH02-014
% titre: Je n'ai pas de talents particuliers
% chapitre_id: 2
% chapitre_nom: Nombres relatifs
% annee_scolaire: 9H
% difficulte: 2
% tags: [entiers-relatifs, addition, soustraction, multiplication, division]
% auteur: Bussy D.
% ---

\begin{exercice}{EX-CH02-014}{Je n'ai pas de talents particuliers}

Calcule chaque opération. Classe ensuite les résultats par ordre croissant et reporte les fragments de phrase correspondants pour reconstituer la citation
d'Albert Einstein.

\def\arraystretch{1.25}
\setlength{\tabcolsep}{4pt}
\vspace{2mm}

\begin{center}
\begin{tabular}{c@{\hspace{8mm}}c}

  \begin{tabular}{|>{\centering\arraybackslash}p{6mm}|p{3.5cm}|p{3cm}|}
    \hline
    \cellcolor{gray!40}\textbf{1.}  &Je& $(+25)+(-1)=$               \\ \hline
    \cellcolor{gray!40}\textbf{2.}  &Je suis& $(-3)\cdot(+7)=$              \\ \hline
    \cellcolor{gray!40}\textbf{3.}  &donné& $(-70)\div(-7)=$              \\ \hline
    \cellcolor{gray!40}\textbf{4.}  &bon& $(+2)\cdot(-1)=$               \\ \hline
    \cellcolor{gray!40}\textbf{5.}  &et& $(-6)\cdot(+12)=$             \\ \hline
    \cellcolor{gray!40}\textbf{6.}  &Dieu& $(+16)\cdot(-4)=$             \\ \hline
    \cellcolor{gray!40}\textbf{7.}  &de& $(+8)-(-5)=$        \\ \hline
    \cellcolor{gray!40}\textbf{8.}  &un& $(-72)+(-8)=$                 \\ \hline
    \cellcolor{gray!40}\textbf{9.}  &juste& $(-21)+(-4)=$                 \\ \hline
    \cellcolor{gray!40}\textbf{10.} &d'un& $(-55)\div(+5)=$              \\ \hline
    \cellcolor{gray!40}\textbf{11.} &De plus,& $(+12)-(-4)=$      \\ \hline
  \end{tabular}

  &

  \begin{tabular}{|>{\centering\arraybackslash}p{6mm}|p{3.5cm}|p{3cm}|}
    \hline
    \cellcolor{gray!40}\textbf{12.} &n'ai& $(+24)\cdot(-2)=$             \\ \hline
    \cellcolor{gray!40}\textbf{13.} &pas& $(-48)\div(-6)=$               \\ \hline
    \cellcolor{gray!40}\textbf{14.} &l'entêtement& $(+10)+(-65)=$                \\ \hline
    \cellcolor{gray!40}\textbf{15.} &curieux& $(-5)+(+17)=$                 \\ \hline
    \cellcolor{gray!40}\textbf{16.} &particulier.& $(-9)\div(+3)=$                \\ \hline
    \cellcolor{gray!40}\textbf{17.} &mulet& $(-11)-(-5)=$       \\ \hline
    \cellcolor{gray!40}\textbf{18.} &assez& $(-80)\div(-40)=$              \\ \hline
    \cellcolor{gray!40}\textbf{19.} &m'a& $(-64)-(+6)=$      \\ \hline
    \cellcolor{gray!40}\textbf{20.} &passionnément& $(-25)\div(+5)=$               \\ \hline
    \cellcolor{gray!40}\textbf{21.} &flair& $(-2)+(+27)=$                 \\ \hline
    \cellcolor{gray!40}\textbf{22.} &talent& $(+13)-(+22)=$     \\ \hline
  \end{tabular}

\end{tabular}
\end{center}

\bigskip
\noindent Je n'ai pas de talents \dotfill \\
\vspace{1cm}
\noindent\dotfill\\[0.5cm]
\dotfill
\makebox[\textwidth]{\dotfill\ « Albert Einstein »}

\end{exercice}


\begin{solution}{EX-CH02-014}

\renewcommand{\arraystretch}{1.7}
\setlength{\tabcolsep}{4pt}
\vspace{2mm}

\begin{center}
\begin{tabular}{c@{\hspace{8mm}}c}

  \begin{tabular}{|>{\centering\arraybackslash}p{6mm}|p{2.5cm}|p{4.0cm}|}
    \hline
    \cellcolor{gray!40}\textbf{1.}  &Je& $(+25)+(-1)=\rc{+24}$               \\ \hline
    \cellcolor{gray!40}\textbf{2.}  &Je suis& $(-3)\cdot(+7)=\rc{-21}$              \\ \hline
    \cellcolor{gray!40}\textbf{3.}  &donné& $(-70)\div(-7)=\rc{+10}$              \\ \hline
    \cellcolor{gray!40}\textbf{4.}  &bon& $(+2)\cdot(-1)=\rc{-2}$               \\ \hline
    \cellcolor{gray!40}\textbf{5.}  &et& $(-6)\cdot(+12)=\rc{-72}$             \\ \hline
    \cellcolor{gray!40}\textbf{6.}  &Dieu& $(+16)\cdot(-4)=\rc{-64}$             \\ \hline
    \cellcolor{gray!40}\textbf{7.}  &de& $(+8)-(-5)=\rc{(+8)+(+5)=+13}$        \\ \hline
    \cellcolor{gray!40}\textbf{8.}  &un& $(-72)+(-8)=\rc{-80}$                 \\ \hline
    \cellcolor{gray!40}\textbf{9.}  &juste& $(-21)+(-4)=\rc{-25}$                 \\ \hline
    \cellcolor{gray!40}\textbf{10.} &d'un& $(-55)\div(+5)=\rc{-11}$              \\ \hline
    \cellcolor{gray!40}\textbf{11.} &De plus,& $(+12)-(-4)=\rc{(+12)+(+4)=+16}$      \\ \hline
  \end{tabular}

  &

  \begin{tabular}{|>{\centering\arraybackslash}p{6mm}|p{2.5cm}|p{4.0cm}|}
    \hline
    \cellcolor{gray!40}\textbf{12.} &n'ai& $(+24)\cdot(-2)=\rc{-48}$             \\ \hline
    \cellcolor{gray!40}\textbf{13.} &pas& $(-48)\div(-6)=\rc{+8}$               \\ \hline
    \cellcolor{gray!40}\textbf{14.} &l'entêtement& $(+10)+(-65)=\rc{-55}$                \\ \hline
    \cellcolor{gray!40}\textbf{15.} &curieux& $(-5)+(+17)=\rc{+12}$                 \\ \hline
    \cellcolor{gray!40}\textbf{16.} &particulier.& $(-9)\div(+3)=\rc{-3}$                \\ \hline
    \cellcolor{gray!40}\textbf{17.} &mulet& $(-11)-(-5)=\rc{(-11)+(+5)=-6}$       \\ \hline
    \cellcolor{gray!40}\textbf{18.} &assez& $(-80)\div(-40)=\rc{+2}$              \\ \hline
    \cellcolor{gray!40}\textbf{19.} &m'a& $(-64)-(+6)=\rc{(-64)+(-6)=-70}$      \\ \hline
    \cellcolor{gray!40}\textbf{20.} &passionnément& $(-25)\div(+5)=\rc{-5}$               \\ \hline
    \cellcolor{gray!40}\textbf{21.} &flair& $(-2)+(+27)=\rc{+25}$                 \\ \hline
    \cellcolor{gray!40}\textbf{22.} &talent& $(+13)-(+22)=\rc{(+13)+(-22)=-9}$     \\ \hline
  \end{tabular}

\end{tabular}
\end{center}

\bigskip
\noindent Ordre croissant des résultats :
\[
  1 - 12 - 13 - 7 - 22 - 16 - 2 - 9 - 20 - 15 - 11 - 6 - 19 - 3 - 14 - 10 - 17 - 5 - 8 - 18 - 4 - 21
\]

\bigskip
\noindent Citation reconstituée :\\
Je n'ai pas de talents \rc{particuliers. Je suis juste passionnément curieux.
De plus Dieu m'a donné l'entêtement d'un mulet, et un assez bon flair.}

\makebox[\textwidth]{\hfill « Albert Einstein »}

\end{solution}